%% LyX 2.5.1 created this file.  For more info, see https://www.lyx.org/.
%% Do not edit unless you really know what you are doing.
\documentclass[english,brazil]{article}
\usepackage[T1]{fontenc}
\usepackage[utf8]{inputenc}
\usepackage{babel}
\usepackage{array}
\usepackage{multirow}
\usepackage{footnotehyper}
\usepackage{amsmath}
\usepackage{accents}
\usepackage{geometry}
\geometry{verbose,tmargin=3cm,bmargin=3cm,lmargin=3cm,rmargin=2.5cm}
\usepackage[]
 {hyperref}

\makeatletter

%%%%%%%%%%%%%%%%%%%%%%%%%%%%%% LyX specific LaTeX commands.

\makesavenoteenv{tabular}

\newcommand{\lyxmathsym}[1]{\ifmmode\begingroup\def\b@ld{bold}
  \text{\ifx\math@version\b@ld\bfseries\fi#1}\endgroup\else#1\fi}

\newcommand{\docedilla}[2]{\underaccent{#1\mathchar'30}{#2}}
\newcommand{\cedilla}[1]{\mathpalette\docedilla{#1}}

%% Because html converters don't know tabularnewline
\providecommand{\tabularnewline}{\\}

\makeatother

\begin{document}
\title{Contabilidade Social}
\maketitle

\subsection{Conceitos básicos}

O princípio das partidas dobradas, reza que, a um lançamento a débito, deve sempre corresponder um outro de mesmo valor a crédito. O equilíbrio interno refere-se à exigência de igualdade entre o valor do débito e o do crédito em cada uma das contas, enquanto o equilíbrio externo implica a necessidade de equilíbrio entre todas as contas do sistema.

Escolhida a contabilidade como o instrumento de aferição macroscópica do movimento econômico , tudo se passa, então, como se a economia de todo um país pudesse ser vista como a de uma única grande empresa: os resultados de seu funcionamento durante um determinado período de temposão apresentados pelas contas integrantes do sistema de contas nacionais.

O que se convencionou chamar “contabilidade social” não se reduz ao sistema de contas nacionais, mas inclui outras peças chave como o balanço de pagamentos, as contas do sistema monetário e os indicadores sociais, como distribuição de renda e índice de desenvolvimento humano. Por isso, a analogia deve resumir-se tão somente à forma -- como contas, balancetes e lançamentos contábeis. Sua substância e seus objetivos são inteiramente distintos.

Na medida em que a contabilidade de que estamos falando é social, toda a “trabalheira” estatística tem de servir para que a sociedade civil possa ter uma ideia mais clara dos rumos de um país e possam, assim, intervir nesses rumos, quando for o caso. 

No sistema econômico em que vivemos, tudo pode ser avaliado monetariamente. Assim, a imensa gama de diferentes bens e serviços que uma economia é capaz de produzir pode ser transformada numa coisa de mesma substância, ou seja, dinheiro. É isso que torna possível a mensuração dos agregados como o produto nacional e a renda nacional.

Uma das noções fundamentais da contabilidade social é a de identidade. Mas aqui identidade parece ter um significado particular, é uma identidade contável, por isso costuma ser utilizado o símbolo da identidade $\equiv$ e não o da igualdade $=$ . Se pegarmos a troca como exemplo, ela implica em duas operações, que são o inverso uma da outra -- o comprador troca \$ 10 por uma camisa e o vendedor troca uma camisa por \$ 10 --, mas que, do ponto de vista analítico, conformam uma identidade, já que uma não pode existir sem a outra. Exemplos: da mesma forma que não pode ocorrer uma compra sem venda, também não pode haver uma produção que não constitua um dispêndio e não seja simultaneamente geração de renda. É a partir da identidade produto  $\equiv$ dispêndio $\equiv$ renda que se deriva o fluxo circular da renda. . Similarmente, poupança implica necessariamente investimento, e investimento não pode ser entendido sem que o consideremos, em contrapartida, como poupança. 

Para se chegar ao produto agregado da economia é preciso deduzir, do valor bruto da produção, o valor do consumo intermediário, isto é, do que foi utilizado como insumo na produção. Aqui por bem final devemos entender um bem é vendido para os consumidors finais, isto é, que não vai ser utilizado como insumo. Por definição, todo bem final deve ter seu valor considerado no cálculo do valor do produto, mas nem todo bem cujo valor entra no cálculo do produto é um bem final por natureza. Pois aqueles produtos que não foram consumidos ainda, e não são considerados produtos finais (isto é, foram produzidos para sere utilizados como insumos) também entram no ponto final. Isto é, o produto agregado da economia é a combinação dos produtos vendidos aos consumidores finais e os produtos não consumidos no período. Tendo em conta que consumidor final é o agente econômico que adquire um bem ou serviço que não será utilizado para revendê-lo ou empregá-lo como insumo em um novo processo produtivo.

Essa forma de abordar é conhecida como ótica da despesa e avalia o produto de uma economia considerando a soma dos valores de todos os bens e serviços produzidos no período que não foram destruídos, ou absorvidos como insumos, na produção de outros bens e serviços. Porém ainda há mais duas formas de abordar a questão. Pela ótica do produto, a avaliação do produto total da economia consiste na consideração do valor efetivamente adicionado pelo processo de produção em cada unidade produtiva. Isto é, se olharmos empresa a empresa, o valor adicionado é dado pela diferença entre o valor de sua produção e o valor da produção que ele adquiriu como insumo. Esta é efetivamente, a contribuição da empresa para a constituição do produto total da economia

Ambos os resultados devem ser idênticos. Mas ainda há outra forma de abordar é pela ótica da renda. Precisamos começar entendendo aquilo que se chama de processo de produção: consideremos a princípio que só existam dois fatores de produção: trabalho e um outro a que daremos genericamente o nome de capital -- concretamente isso envolve além das máquinas e demais equipamentos, todo o conjunto de elementos que conformam as condições objetivas sem as quais o processo de produção não pode acontecer. É entre capital e trabalho, portanto, que deve ser repartido o produto gerado pela economia, pois foi sua participação no processo produtivo que garantiu a obtenção desse produto\footnote{É possível além destes dois fatores a terra, mas pouco muda a discussão.}. À remuneração do fator trabalho damos o nome de salário e à remuneração do fator capital damos o nome de lucro\footnote{E a terra seria aluguel.}.

Assim, num dado período de tempo, as remunerações de ambos os fatores, conjuntamente consideradas, devem igualar, em valor, o produto obtido pela economia nesse mesmo período, visto não ser as remunerações nada mais do que a divisão do produto. As remunerações pagas constituem o que chamamos de renda. Não é preciso muito esforço para perceber que, com isso, consuma se a identidade produto $\equiv$ renda. Logo, o produto gerado por uma economia num determinado período de tempo é igual à renda gerada nesse mesmo período.

Algumas observações extras, é que aqui, diferente do marxismo, dinheiro e moeda são apresentados como sinônimos, valor e preço também. Além disso a sociedade é dividida em dois grupos de forma diferente. Lê-se: ``Os membros que constituem a sociedade aparecem duas vezes no jogo de sua reprodução material e desempenham dois papéis distintos: num determinado momento, são produtores; no outro, surgem como consumidores daquilo que foi produzido''. Ou seja, o trabalhador e o capitalista são ambos produtores e ambos consumidores. Isso também fica claro evidente quando vemos nos fatores de produção o ``trabalho'' e o ``capital'' lado a lado,cada um sendo responsável por produzir ``valor'' igualmente. Ou seja, no processo de produção o trabalhador contribui com trabalho, e o capitalista com capital. Ambos em iramndade. Isso também aparece no próximo destaque que o livro dá: Além de desempenhar o papel de consumidores, as famílias detêm também a condição de proprietárias dos fatores de produção e é nessa condição que elas garantem seu acesso aos bens e serviços produzidos. Aqui, por ``propriedade dos fatores de produção'', temos que os trabalhadores só tem a força de trabalho para vender.

Temos então que a ótica do produto considera a atividade dos indivíduos como produtores, ou seja, a atividade das unidades produtivas ou empresas. Finalmente, a ótica da renda analisa os indivíduos em sua condição de proprietários de fatores de produção. As transações ocorrem entre famílias e empresas e envolvem fluxos reciprocamente determinados de bens e serviços concretos, por um lado, e de dinheiro, por outro. Seguindo o livro, temos este fluxo:
\begin{enumerate}
\item As famílias cedem às empresas os fatores de produção de que são proprietárias e, em troca, recebem das empresas uma renda, ou seja, uma remuneração sob a forma de dinheiro;
\item As empresas combinam esses fatores num processo denominado processo de produção e obtêm, como resultado, um conjunto de bens e serviços;
\item Com a renda recebida em troca da utilização, na produção, dos fatores de que são proprietárias, as famílias compram das empresas os bens e serviços por estas produzidos;
\item As famílias consomem os bens e serviços;
\end{enumerate}
Observe como esse fluxo coloca em igualdade 'famílias' de trabalhadores e famílias de capitalistas. A renda (salário ou lucro) é visto como uma ``gratificação'' pela sua contribuição no processo produtivo. Concluindo o capítulo, precisamos ter claramente que o fluxo é da renda e não do produto: o dinheiro que remunera os fatores de produção é o mesmo que retorna às empresas na compra dos bens e serviços finais. Isso não acontece com os demais bens.

\subsection{Estrutura básica das contas nacionais}

As maiores simplificações normalmente são considerar ecnomia fechada e sem governo, vamos trabalhar com este exemplo.

Ora, tudo aquilo que é produzido em um período, mas que não é consumido nesse período, significando, ou ensejando, consumo no futuro, tem um nome: chama-se investimento. O investimento divide-se em formação de capital fixo e variação de estoques. O investimento costuma ser dividido em variação de estoques, que congrega os bens cujo consumo futuros irão se dar de uma única vez, e a formação bruta de capital fixo, que agrega os bens que não desaparecem depois de uma única utilização e possibilitam a produção ao longo de um determinado período de tempo.

O desgaste do capital fixo chama-se depreciação: Os bens considerados capital fixo também se desgastam com o tempo e com o uso, de modo que, findo um determinado período, seu valor terá sido inteiramente absorvido pelo fluxo de produção de bens aí ocorrido. Para obter o valor do produto líquido de uma economia em um determinado período é preciso deduzir, do valor do produto bruto, a parcela destinada à reposição do estoque de capital da economia, ou seja, a depreciação.
\begin{center}
\begin{tabular}{|c|c|c|c|}
\hline 
 & \textbf{Débito} &  & \textbf{Crédito}\tabularnewline
\hline 
\textbf{A} & Produto Líquido & \textbf{C} & Consumo Pessoal\tabularnewline
\hline 
\textbf{B} & Depreciação & \textbf{D} & Variação de estoques\tabularnewline
\hline 
 &  & \textbf{E} & Formação de capital fixo\tabularnewline
\hline 
 & \textbf{Produto Bruto} &  & \textbf{Despesa Bruta}\tabularnewline
\hline 
\end{tabular}
\par\end{center}

Quanto à remunerações, além de salários e lucros, precisamos introduzir duas outras categorias de remuneração: os aluguéis (que remuneram os proprietários de imóveis de modo geral) e os juros (que remuneram os proprietários de capital monetário). O único cuidado adicional que deve ser tomado é evitar a dupla contagem que pode ocorrer se considerarmos nessas rubricas, além dos aluguéis e dos juros pagos às famílias, também aqueles pagos às empresas. Estes últimos não devem ser considerados porque, como receitas, já participam dos demonstrativos de lucros e perdas das empresas. Com exceção do setor finnceiro, mas isso será visto adiante.
\begin{center}
\begin{tabular}{|c|c|c|c|}
\hline 
 & \textbf{Débito} &  & \textbf{Crédito}\tabularnewline
\hline 
\hline 
$\boldsymbol{a}_{1}$ & salários & \textbf{C} & Consumo Pessoal\tabularnewline
\hline 
$\boldsymbol{a}_{2}$ & lucros & \textbf{D} & Variação de estoques\tabularnewline
\hline 
$\boldsymbol{a}_{3}$ & alugueis & \textbf{E} & Formação de capital fixo\tabularnewline
\hline 
$\boldsymbol{a}_{4}$ & juros &  & \tabularnewline
\hline 
\textbf{A} & Produto líquido\footnote{Ou renda ou produto nacional líquido.} - $A=\sum a_{i}$ &  & \tabularnewline
\hline 
\textbf{B} & Depreciação &  & \tabularnewline
\hline 
 & \textbf{Produto Bruto}{\footnote{Ou Renda.}} &  & \textbf{Despesa Bruta}\tabularnewline
\hline 
\end{tabular}
\par\end{center}

A conta de produção mostra a identidade entre renda e dispêndio, enquanto a conta de apropriação mostra de que maneira as famílias alocam as rendas recebidas pela cessão deseus fatores de produção às empresas.

Já a conta de capital mostra a identidade investimento $\equiv$ poupança, que nada mais é do que uma forma alternativa de representar a identidade produto $\equiv$ renda $\equiv$ dispêndio.

\subsubsection{Abrindo a economia}

Até agora trabalhamos com o caso mais simples e que eu devo focar minha pesquisa, mas é interessante relaxarmos um pouco as hipóteses. Para isso precisamos ter mais alguns conceitos.
\begin{itemize}
\item \textbf{Balança comercial}: o comércio de bens com o resto do mundo.
\item \textbf{Balança de serviços e rendas:} comércio de serviços e fatores de produção (rendas como lucros e juros).
\end{itemize}
Agora temos então uma distinção entre produto interno e produto nacional. Não há uma resposta, mas de acordo com o SNA 93 usamos o atributo ``interno'' para ``produto'' (produto interno), e ``nacional'' para ``renda nacional''. A ideia é que normalmennte temos interesse por um lado no PIB (produto interno bruto), isto é, o produto total o produto total produzido no território do país, independentemente da origem dos fatores de produção responsáveis por ele. E de outro lado, a RNB (renda nacional bruta) no valor adicionado gerado por fatores de produção de propriedade de residentes, independentemente do território onde esse valor é gerado.

Para se obter o produto nacional de uma economia, é preciso deduzir de seu produto interno a renda líquida enviada ao exterior ou, se for o caso, adicionar a seu produto interno a renda líquida recebida do exterior. Na maior parte dos casos, os países mais desenvolvidos são exportadores líquidos de capital e, portanto, recebem liquidamente renda do exterior, enquanto ocorre o inverso com os países menos desenvolvidos. Vamos agora para a terceira versão da conta do produto:
\begin{center}
\begin{tabular}{|c|c|c|c|}
\hline 
 & \textbf{Débito} &  & \textbf{Crédito}\tabularnewline
\hline 
\hline 
\textbf{I} & Importação de bens e serviços & \textbf{G} & Exportações de bens e serviços\tabularnewline
\hline 
\textbf{J-H} & Renda líquida enviada {[}+{]} ao exterior & \textbf{C} & Consumo Pessoal\tabularnewline
\hline 
$\boldsymbol{a}_{1}$ & salários & \textbf{D} & Variação de estoques\tabularnewline
\hline 
$\boldsymbol{a}_{2}$ & lucros & \textbf{E} & Formação de capital fixo\tabularnewline
\hline 
$\boldsymbol{a}_{3}$ & alugueis &  & \tabularnewline
\hline 
$\boldsymbol{a}_{4}$ & juros &  & \tabularnewline
\hline 
\textbf{A} & Produto líquido\footnote{Ou renda ou produto nacional líquido.} ( $A=\sum a_{i}$) &  & \tabularnewline
\hline 
\textbf{B} & Depreciação &  & \tabularnewline
\hline 
 & \textbf{Oferta de Bens e Serviços} &  & \textbf{Demanda por Bens e Serviços}\tabularnewline
\hline 
\end{tabular}
\par\end{center}

Devemos notar que \textbf{G} e \textbf{I} representam as exportações e importações de bens e serviços não fatores, compondo a balança comercial e parcialmente a balança de serviços e rendas. Já \textbf{H} e \textbf{J} representam as rendas recebidas do exterior e enviadas ao exterior, correspondendo ao restante balança de serviços e rendas. 

Pode-se dizer que a conta de produção apresenta, do lado do débito, o PIB mais as importações de bens e serviços não fatores, ou seja, a oferta total, que deve igualar em valor à demanda total por bens e serviços. Mas veja que estamos impondo esta igualdade entre oferta e demanda por por definições. Ms é evidente que não podemos contar duas vezes, então uma versão mais sucinta pode ser dada apenas por:
\begin{center}
\begin{tabular}{|c|c|c|c|}
\hline 
 & \textbf{Débito} &  & \textbf{Crédito}\tabularnewline
\hline 
\hline 
\textbf{I} & Importação de bens e serviços & \textbf{G} & Exportações de bens e serviços\tabularnewline
\hline 
\textbf{J-H} & Renda líquida enviada {[}+{]} ao exterior & \textbf{C} & Consumo Pessoal\tabularnewline
\hline 
\textbf{A} & Renda Líquida ($\sum a_{i}$) & \textbf{D} & Variação de estoques\tabularnewline
\hline 
\textbf{B} & Depreciação & \textbf{E} & Formação de capital fixo\tabularnewline
\hline 
 & \textbf{Oferta de Bens e Serviços} &  & \textbf{Demanda por Bens e Serviços}\tabularnewline
\hline 
\end{tabular}
\par\end{center}

Onde a renda líquida é o agregado dos salários, lucros, alugueis e juros. É importante lembrar que depreciação não é nenhuma transação, apenas a desvalorização, ou melhor diznedo, uma perda de valor do estoque de capital, registrada contabilmente.

O lado do débito reúne as fontes da oferta de bens e serviços disponíveis na economia, já o lado do crédito registra os usos dessa oferta, isto é, para onde os bens e serviços são destinados: consumo, investimento em estoques, formação de capital fixo e exportações. Assim, débito corresponde à origem ou oferta dos recursos, enquanto crédito corresponde ao seu destino.

\subsubsection{Colocando governo}

O governo arrecada impostos diretos (que incidem sobre a renda ou o patrimônio, como IPTU, IPVA, etc) e indiretos (que incidem sobre os preços, como ICMS, IPI, etc). Transferências são impostos diretos com o sinal negativo; subsídios são impostos indiretos com o sinal negativo. As duas categorias mais importantes de transferência são, por um lado, as pensões e aposentadorias e, por outro, os juros da dívida pública.

Tributo é a designação genérica de todo tipo de renda que o governo é capaz de arrecadar por ser governo. Os impostos (diretos e indiretos) são os tributos mais conhecidos, mas, além deles, existem também as taxas (como a taxa do lixo), as contribuições de melhoria (decorrentes da realização de obras públicas) e outros tipos de contribuição (como a previdenciária, a de intervenção no domínio econômico etc.).

Quanto aos subsídios, na maior parte das vezes, os subsídios não significam propriamente a redistribuição de uma receita coletada por meio de impostos, mas simplesmente a abdicação, por parte do governo, de uma receita à qual ele teria direito. O governo pode, por exemplo, querer reduzir o preço do leite aos consumidores finais e para tanto, abrir mão da arrecadação do imposto sobre circulação de mercadorias e serviços (ICMS) que incidiria sobre a comercialização do leite.

Para lidar com a modificaçaõ dos preços devidos a impostos e subsídios, foram criados dois conceitos de produto: o produto a preços de mercado, que inclui o valor dos impostos indiretos compensados dos subsídios, e o produto a custo de fatores , que não considera esse valor adicional.

Agora oferta global da economia que discutíamos previamente em um determinado período é a soma do produto interno bruto a preços de mercado com as importações de bens e serviços não fatores. A demanda global, por seu lado, é a soma do PIB também a preços de mercado com as exportações de bens e serviços não fatores. Agora então o quadro de produção é: 
\begin{center}
\begin{tabular}{|c|c|c|c|}
\hline 
 & \textbf{Débito} &  & \textbf{Crédito}\tabularnewline
\hline 
\hline 
\textbf{I} & Importação de bens e serviços & \textbf{G} & Exportações de bens e serviços\tabularnewline
\hline 
\textbf{J-H} & Renda líquida enviada {[}+{]} ao exterior & \textbf{C} & Consumo Pessoal\tabularnewline
\hline 
\textbf{A} & Renda Líquida ($\sum a_{i}$) & \textbf{L} & Consumo do governo\tabularnewline
\hline 
\textbf{B} & Depreciação & \textbf{D} & Variação de estoques\tabularnewline
\hline 
\textbf{Q-N} & Impostos indiretos líquidos\footnote{Isto é, os impostos indiretos \textbf{Q} descontado dos subsídios \textbf{N}.} & \textbf{E} & Formação de capital fixo\tabularnewline
\hline 
 & \textbf{Oferta de Bens e Serviços} &  & \textbf{Demanda por Bens e Serviços}\tabularnewline
\hline 
\end{tabular}
\par\end{center}

Notamos que agora o governo surge como uma gente adicional de demanda. Mas que não encontramos diretamente as traansferêncis ou impostos diretos, nem o saldo do governo em conta corrent e outras receitas correntes líquidas, pois aqui estamos olhando para o produto, portanto olhamos apenas o que afeta diretamente o produto. IEstes outros elementos devem aparecer por exemplo, na conta de apropriação:
\begin{center}
\begin{tabular}{|c|c|c|c|c|c|}
\hline 
 & \textbf{Débito} &  & \multicolumn{3}{c|}{\textbf{Crédito}}\tabularnewline
\hline 
\hline 
\textbf{C} & Consumo Pessoal & $\boldsymbol{a}_{1}$ & salários & \multirow{4}{*}{\textbf{A}} & \multirow{4}{*}{Renda Líquida}\tabularnewline
\cline{1-4}
\textbf{P-M} & Impostos diretos líquidos\footnote{Isto é, os impostos diretos \textbf{P} descontado das transferências \textbf{M}.} & $\boldsymbol{a}_{2}$ & lucros &  & \tabularnewline
\cline{1-4}
\textbf{R} & Outras receitas corrents líquidas & $\boldsymbol{a}_{3}$ & alugueis &  & \tabularnewline
\cline{1-4}
\textbf{F} & Poupança privada líquida & $\boldsymbol{a}_{4}$ & juros &  & \tabularnewline
\hline 
 & \textbf{Utilização da renda Líquida} &  & \textbf{Renda líquida} &  & \tabularnewline
\hline 
\end{tabular}
\par\end{center}

E para entender a poupança privada líquida, olhamos a conta de capital:
\begin{center}
\begin{tabular}{|c|c|c|c|}
\hline 
 & \textbf{Débito} &  & \textbf{Crédito}\tabularnewline
\hline 
\hline 
\textbf{D} & Variação de estoques & \textbf{F} & Poupança privada líquida\tabularnewline
\hline 
\textbf{E} & Formação bruta de capital fixo & \textbf{B} & Depreciação\tabularnewline
\hline 
 &  & \textbf{K} & Resultado do BP em transações corrente\tabularnewline
\hline 
 &  & \textbf{O} & Saldo do governo em conta corrente\tabularnewline
\hline 
 & \textbf{Investimento Bruto Total} &  & \textbf{Proupança Bruta Total}\tabularnewline
\hline 
\end{tabular}
\par\end{center}

Se fossemos retomar uma economia fechada sem governo, em crédito teriamos apenas \textbf{F} e \textbf{B}. Ou seja, a a poupança privada líquida é a produção que não foi consumida descontada a depreciação, ou seja é \textbf{F=D+E-B, }onde B é a depreciação, D a variação de estoques e E a formação bruta de capital fixo. Outras tabelas que eu não vou elaborar ainda são a conta do setor externo:
\begin{center}
\begin{tabular}{|c|c|c|c|}
\hline 
 & \textbf{Débito} &  & \textbf{Crédito}\tabularnewline
\hline 
\hline 
\textbf{G} & Exportações de bens e serviços & \textbf{I} & Importação de bens e serviços\tabularnewline
\hline 
\textbf{H} & Renda recebida do exterior & \textbf{J} & Renda enviada ao exterior\tabularnewline
\hline 
\textbf{K} & Resultado do BP em transações corrents &  & \tabularnewline
\hline 
 & \textbf{Total do Débito} &  & \textbf{Total do Crédito}\tabularnewline
\hline 
\end{tabular}
\par\end{center}

E a conta do governo. 
\begin{center}
\begin{tabular}{|c|c|c|c|}
\hline 
 & \textbf{Débito} &  & \textbf{Crédito}\tabularnewline
\hline 
\hline 
\textbf{L} & Consumo do governo & \textbf{P} & Impostos diretos\footnote{Para empresas ($p.1$) e famílias ($p.2$).}.\tabularnewline
\hline 
\textbf{M} & Transferências & \textbf{Q} & Impostos indiretos\tabularnewline
\hline 
\textbf{N} & Subsídios & \textbf{R} & Outras receitas correntes líqudas\tabularnewline
\hline 
\textbf{O} & Saldo do governo em conta corrente &  & \tabularnewline
\hline 
 & \textbf{Utilização da Receita} &  & \textbf{Total da Receita}\tabularnewline
\hline 
\end{tabular}
\par\end{center}

A rubrica \textbf{R} também tem de ser lançada uma vez que os detentores de renda também fazem diretamente, além dos impostos diretos propriamente ditos, outros tipos de pagamento ao governo (como o pagamento de taxas, contribuições previdenciárias, multas ou aluguéis em função do uso de propriedades do Estado), os quais vêm a constituir as outras receitas do governo.

E agora também precisamos considerar \textbf{O}, pois uma vez introduzido no sistema, o governo torna-se também uma fonte geradora de poupança, assim sendo a terceira fonte geradora de poupança (setor privado, setor externo e governo).

O equilíbrio interno das contas está supostamente garantido. Uma forma de conferirmos isso é somarmos o lado do débito de todas as cinco contas e deduzirmos disso o somatório do lado do crédito de todas as contas. Se o sistema de fato estiver equilibrado o resultado dessa operação deverá ser zero.
\begin{center}
{[}(I+J-H+A+B+Q-N)+(C+P-M+R+F)+(D+E)+(G+H+K+L+M+N+O){]}
\par\end{center}

\begin{center}
-
\par\end{center}

\begin{center}
{[}(G+C+L+D+E)+(A)+(F+B+K+O)+(I+J)+(P+Q+R){]}
\par\end{center}

\begin{center}
= 0
\par\end{center}

É importante destacar aqui que este é um modelo de contas nacionais (SNA 93) apenas, que não é o que está em uso atualmente, mas por ser simpificado eu considerei mais interessante explorar. Além disso eu dei ênfase na conta do produto. 


\subsubsection{Matriz insumo-produto}

A matriz insumo-produto, cujo desenvolvimento está ligado ao prêmio Nobel em Economia Wassily W. Leontief (1906-1999), tem como objetivo proporcionar uma análise acerca das relações intersetoriais na produção, podendo ser vista como um instrumento alternativo ao sistema de contas nacionais.

A partir de uma matriz insumo-produto, pode-se, por exemplo, estimar qual é o impacto sobre o nível de produção e emprego e sobre as demandas setoriais, de um aumento ou uma retração na produção de um determinado ramo (um tipo de informação que um sistema convencional de contas nacionais não é capaz de fornecer).

Um exemplo bem simples pode ser útil para compreender a ideia da matriz insumo-produto, bem como sua forma de funcionamento e sua utilidade. Consideremos uma economia hipotética com apenas três setores -- digamos 1, 2 e 3 -- que estabelecem transações econômicas entre si. Se $X_{ij}$ representa as vendas do setor i para o setor j, podemos construir a matriz:
\begin{center}
\begin{tabular}{|c|c|c|c|c|c|c|}
\hline 
 &  & \multicolumn{3}{c|}{Setores} & Demanda & Produção\tabularnewline
\hline 
\hline 
 &  & 1 & 2 & 3 &  & \tabularnewline
\hline 
\multirow{3}{*}{Setores} & 1 & $x_{11}$ & $x_{12}$ & $x_{13}$ & $y_{1}$ & $x_{1}$\tabularnewline
\cline{2-7}
 & 2 & $x_{21}$ & $x_{22}$ & $x_{33}$ & $y_{2}$ & $x_{2}$\tabularnewline
\cline{2-7}
 & 3 & $x_{31}$ & $x_{32}$ & $x_{33}$ & $y_{3}$ & $x_{3}$\tabularnewline
\hline 
Valor adicionado &  & $v_{1}$ & $v_{1}$ & $v_{3}$ &  & \tabularnewline
\hline 
Produção &  & $x_{1}$ & $x_{2}$ & $x_{3}$ &  & \tabularnewline
\hline 
\end{tabular}
\par\end{center}

Considerando as vendas do setor i para o setor j como uam fraçao da produção total do setor j então:

\[
x_{ij}=a_{ij}x_{j}
\]

Então podemos ter uma matriz de oeficientes técnicos:
\begin{center}
\begin{tabular}{|c|c|c|c|}
\hline 
 & 1 & 2 & 3\tabularnewline
\hline 
1 & $a_{11}$ & $a_{12}$ & $a_{13}$\tabularnewline
\hline 
2 & $a_{21}$ & $a_{22}$ & $a_{23}$\tabularnewline
\hline 
3 & $a_{31}$ & $a_{32}$ & $a_{33}$\tabularnewline
\hline 
\end{tabular}
\par\end{center}

Então podemos consruir:

\[
AX+Y=X
\]

Onde $A$ é a mariz de coeficientes técnicos, e $X$($Y$) a matriz de uma coluna dos elementos $x_{i}$($y_{i}$). Ou seja, podemos pensar de outra forma, qual é a demanda pelo setor $i$? Vai ser vai ser a ``produção'' líquida de $x_{i}$, isso é, $x_{i}$ descontado as vendas que o setor $i$ precisou fazer para todos os outros setores:

\begin{align*}
y_{i} & =x_{i}-\left(\sum_{j}x_{ij}\right)=x_{i}-\left(\sum_{j}a_{ij}x_{j}\right)
\end{align*}

Ou rearranjando, e no nosso caso detrês setores:

\[
x_{i}=a_{i1}x_{1}+a_{i2}x_{2}+a_{i3}x_{3}+y_{i}
\]
Valo observar que $a_{ii}$ indica a ``venda'' do setor i para o opróprio setor i, podemos supor que é apenas uma parte da produção do setor i que é consumido pelo próprio setor na sua produção. Ainda podemos rearranjar a formulação matricial como:

\begin{align*}
AX+Y & =X\\
\left(I-A\right)X & =Y\\
X & =\left(I-A\right)^{-1}Y\\
X & =LY
\end{align*}

Onde $L$ é a matriz de Leontieff caculda a partir da matriz de coeficientes. Com ela, e com determinada demanda fixa $Y$, podemos calcular então qua deve ser a produção de cada setor para atender esta demanda. Vamos trabalhar então com um exemplo oposto do livro, código pode ser encontrado no \href{https://github.com/jdansb/Econophysics/blob/main/files/contsocial.ipynb}{meu GitHub}. Vamos partir da matriz de coeficientes:
\begin{center}
\begin{tabular}{|c|c|c|c|}
\hline 
 & 1 & 2 & 3\tabularnewline
\hline 
1 & 0.09 & 0.08 & 0.01\tabularnewline
\hline 
2 & 0.18 & 0.21 & 0.11\tabularnewline
\hline 
3 & 0.00 & 0.05 & 0.00\tabularnewline
\hline 
\end{tabular}
\par\end{center}

Então podemos calcular a matriz de Leontieff\footnote{Vamos trabalhar agora com notação vetorial.}:

\[
L=\left(I-A\right)^{-1}=\left[\begin{array}{ccc}
1.12 & 0.11 & 0.02\\
0.26 & 1.30 & 0.15\\
0.01 & 0.06 & 1.01
\end{array}\right]
\]

Se temos então uma demanda de $y=\left[\begin{array}{ccc}
200 & 2000 & 1808\end{array}\right]^{T}$, então a produção bruta total deve ser:

\[
x=Ly=\left[\begin{array}{c}
500\\
2900\\
1952
\end{array}\right]
\]

Obs.: acho que não foi especificado, mas aqui devemos estar medindo tudo em preços. Além disso, sendo $x_{ij}=a_{ij}x_{j}$ podemos recuperar a matriz $X$:
\begin{center}
\begin{tabular}{|c|c|c|c|c|c|c|}
\hline 
 &  & \multicolumn{3}{c|}{Setores} & Demanda & Produção\tabularnewline
\hline 
\hline 
 &  & 1 & 2 & 3 &  & \tabularnewline
\hline 
\multirow{3}{*}{Setores} & 1 & 45 & 240 & 15 & 200 & 500\tabularnewline
\cline{2-7}
 & 2 & 90 & 600 & 210 & 2000 & 2900\tabularnewline
\cline{2-7}
 & 3 & 0 & 144 & 0 & 1808 & 1952\tabularnewline
\hline 
Valor adicionado &  & 365 & 1916 & 1727 & 4008 & \tabularnewline
\hline 
Produção &  & 500 & 2900 & 1952 &  & 5352\tabularnewline
\hline 
\end{tabular}
\par\end{center}

Onde o total da produção total é apenas a soma da produção de cada setor $\sum_{i}x_{i}$. Assim com a demanda total é apenas a soma da demanda de cada setor $\sum_{i}y_{i}$. É uma igualdade contáil que o valor adicionado corresponda a demanda, afinal o valor adicionado total é o valor do excedente. 

Porém para calcular o valor adicionado por setor, precisamos considerar a produção bruta desse setor $x_{i}$ e descontar o que foi vendido de outros setores para ele $\sum_{j}x_{ji}$. Por tanto o valor adicionado por setor é $v_{i}=x_{i}-\sum_{j}x_{ji}$. Por exemplo,o setor $1$ produziu $500$, mas consumiu 45 do próprio setor e 90 do setor 2, então o valor adicionado foi $500-90-45=365$.

Portanto o valor adicionado mede a contribuição de cada setor para a geração de renda, enquanto a demanda final mede o destino final da produção, e coincidem quando agregados para toda a economia.

\subsubsection{Contas nacionais e macroeconomia}

A contabilidade nacional surgiu a partir do advento da teoria keynesiana, quando o economista inglês John Maynard Keynes, em meados dos anos 1930, escreveu a Teoria Geral do Emprego, do Juro e da Moeda para atacar a teoria então vigente e mostrar que a economia não dispunha de mecanismos automáticos para sair de situações de recessão e desemprego. É interessante apontar a abrodagem marginalista do desemprego. Nesta abordagem, todo o desemprego existente era tomado como desemprego voluntário, ou seja, considerava-se que as pessoas que eventualmente não estavam trabalhando encontravam-se em tal situação porque não se dispunham a ofertar sua força de trabalho aos salários vigentes. Em outras palavras, não trabalhavam porque não queriam. A enorme crise dos anos 1930 mostrara a clara inadequabilidade de tal teoria para explicar a realidade. Keynes, portanto, tentou demonstrar que não existia o tal regulador automático e que, por conseguinte, a maior parte do desemprego era involuntário, vale dizer, decorrente de uma demanda por força de trabalho diminuta e, assim, incapaz de empregar toda a oferta existente. Para conseguir demonstrar essa situação, Keynes teve de fazer uma verdadeira revolução nas ideias econômicas e jogar por terra vários dos postulados que constituíam a espinha dorsal da teoria então dominante.

Para Keynes, a teoria econômica com a qual todos os economistas de então trabalhavam estava dominada, implícita ou explicitamente, pela Lei de Say, aquela que diz que toda oferta cria sua própria demanda, decorrendo daí a falta de preocupação com a construção de uma teoria que explicasse os fatores determinantes do comportamento da demanda agregada

Porém o “consenso keynesiano” foi rompido, em meados da década de 1970, pelo advento da teoria das expectativas racionais, que deu nova vida aos pressupostos que Keynes atacara e recuperou a primazia da teoria ortodoxa. A inflação combinada ao desemprego que marcou o final dos anos 1970 levou a uma onda de questionamentos quanto à pertinência da atuação do Estado como regulador do nível de demanda efetiva e, assim, pôs na dianteira as políticas associadas àquilo que se convencionou chamar neoliberalismo (desregulamentação, controle dos gastos públicos, Estadomínimo, privatizações).

É importante frisar que o sistema de contas nacionais pouco ou nada foi abalado por toda essa reviravolta. As identidades macroeconômicas não são, por si só, indicadoras de relações de causalidade entre as variáveis que as constituem.Hoje, a identidade enre produto, renda e dispêndio não é questionada por ninguém independentemente de se aceitar ou não as relações de causa e efeito que as proposições teóricas keynesianas defendem.

Vamos começar então considerando uma economia fechada e sem governo, se tomamos renda por $Y$, consumo por $C$ e investimento por $I$ então: $Y=C+I$. Queremos avançar nossa compreensão, buscando entender na teorya de Keynes, o que determina $C$ e $I$. Segundo Keynes o principal fator a determinar $C$ é $Y$. No entanto, dado um aumento na renda, o aumento do consumo é menos do que proporcional àquele, uma vez que existe aquilo que Keynes chamou propensão a consumir, a qual deriva de algo que ele denominou lei psicológica fundamental. Em outras palavras, dado um determinado nível de renda, as famílias consomem boa parte dela, mas também poupam uma parte. A propensão a consumir é muito maior nas famílias de baixa renda (no limite, as famílias de renda extremamente baixa não poupam nada de sua renda, consumindoa integralmente) e muito maior nas famílias de renda mais elevada. Podemos dizer então que temos uma proprensão $0<c<1$ de consumir, onde $C=f\left(Y\right)$. 

Mas considerando que existe uma parcela do consumo que não varia com o nível de renda ) e que podemos chamar de consumo autônomo, indicado por $C_{a}$, então $C=C_{a}+cY$ e também $Y=C_{a}+cY+I$\footnote{Suponho que se $I=0$, então caso $C_{a}>0$, necessariamente temos $c<1$ para que $Y=C$. E ainda mais explicitamente $C_{a}=\left(1-c\right)Y$.}. Isolando $Y$ temos:

\[
Y=\frac{C_{a}+I}{\left(1-c\right)}
\]

Já quanto ao investimento, Keynes defende que ele depende de duas variáveis. Essas variáveis são a taxa de juros $r$\footnote{Mais especificamente a variável é a preferência pela liquidez, isto é, o desejo das pessoas (e bancos) guardarem dinheiro. Mas segundo o economista inglês, esta preferência está na base da determinação da taxa de juros da economia.} a eficiência marginal do capital $E_{m}$\footnote{Ou mais especificmente as expectativas quanto ao rendimento futuro esperado dos bens de capital, isto é, a expectativa de rentabilidade futura de um novo investimento. Esta expectativa determina então a eficiência marginal.}.Podemos escrever que $I=f\left(r,E_{m}\right)$. Assim o investimento é extremamente instável. Como a teoria keynesiana dos determinantes do investimento é extremamente complexa, explicá-la em detalhes demandaria talvez vários capítulos, o que, com certeza, foge ao escopo deste texto.

Retomando a formula para $Y$, Keynes chamou o termo $1/\left(1-c\right)$ de multiplicador, se escrevermos $M=\left(1-c\right)^{-1}$, então $Y=M\left(C_{a}+I\right)$. Podemos notar que quanto maior a propenção de consumir, isto é, quanto maior $c$, maior o multiplicador $M$. Se $C_{a}$ for estável, então o comportamento de $Y$ fica mais dependente do comportamento de $I$, que é extremamente instável.

Agora vamos considerar a economia aberta e com governo. Além do consumo privado $C$, investimento $I$, vamos considerar os gastos do governo $G$, as exportações líquidas $X$ e as importações $M$.

\[
Y=C+I+G+\left(X-M\right)
\]

12. Os agregados mensurados do ponto de vista interno medem o valor total produzido no território do país, independentemente da origem dos fatores responsáveis por essa produção, enquanto os agregados mensurados do ponto de vista nacional consideram o valor adicionado gerado por fatores de produção de propriedade de residentes, independentemente do território onde esse valor é gerado.

Transpondo para essa expressão ampliada as mesmas considerações anteriormente feitas para uma economia fechada e sem governo, podemos perceber que o nível de renda em que opera a economia não depende apenas do consumo e do investimento, mas também dos gastos do governo e das exportações líquidas das importações. Valem também para essas novas variáveis, as mesmas relações anteriormente estabelecidas para $C_{a}$ e $I$. Então:

\[
Y=\frac{1}{\left(1-c\right)}\left[C_{a}+I+G+\left(X-M\right)\right]
\]

Uma forma bastante sugestiva de compreender esse processo é pensar num mecanismo de estímulos e desestímulos que estão permanentemente influenciando o nível de renda e de produto. Se há um aumento na parcela autônoma do consumo, ou no investimento, ou nos gastos do governo, ou ainda na demanda externa pelos bens e serviços que a economia em questão produz, qualquer um desses aumentos vai estimular a produção e elevar o nível de renda na magnitude determinada pelo multiplicador. 

No caso das exportações, trata-se, na verdade, de um estímulo externo, ou, em outras palavras, de uma injeção de demanda na economia, que provém de um aumento na demanda externa pelos bens e serviços internamente produzidos. Simetricamente, um aumento nas importações representa um vazamento de estímulo, ou seja, uma transferência, para fora da economia, de uma parcela de sua demanda por bens e serviços.

Podemos ainda notar a importância que acabou sendo atribuída ao governo por conta das considerações de Keynes quanto aos determinantes do nível de renda. Se um aumento no nível de renda e produto em que opera a economia pode ser proveniente de uma elevação nos gastos do governo, então cabe a este um importante papel, além daqueles normalmente a ele consagrados.

Uma das óticas de mensuração do produto é justamente a ótica da demanda (dispêndio), a qual mostra a composição do PIB a partir das categorias de demanda tal como estipulado na equação $Y=C+I+G+\left(X-M\right)$.

\subsubsection{Identidades contábeis}

As seguintes identidades derivam das identidades básicas produto$\equiv$ dispêndio$\equiv$ renda e investimento$\equiv$ popupança. Mas antes precisamos ter a distinção entre os agregados mensurados do ponto de vista:
\begin{itemize}
\item \textbf{Interno}: medem o valor total produzido no território do país, independentemente da origem dos fatores responsáveis por essa produção.
\item \textbf{Nacinal}: consideram o valor adicionado gerado por fatores de produção de propriedade de residentes, independentemente do território onde esse valor é gerado. 
\end{itemize}
Então:
\begin{itemize}
\item \textbf{Produto Interno}: reflete o produto total produzido no território do país, independentemente da origem dos fatores de produção responsáveis por ele.
\item \textbf{Renda Naciona}l: considera o valor adicionado gerado por fatores de produção de propriedade de residentes, independentemente do território onde esse valor é gerado.
\end{itemize}
Sendo então:
\begin{itemize}
\item RNB = Renda Nacional bruta
\item PIB = Produto Interno Bruto
\item RLEE = Renda Líquida Enviada ao Exterior
\end{itemize}
Temos 

\[
RNB=PIB-RLEE
\]

Agora se TUR é Transferências Unilaterais Líquidas Recebidas, então:

\[
RDB=RNB+TUR
\]

Onde RDB é Renda nacional Disponível Bruta. Com isso podemos reescrever:

\[
PIB=RDB+RLEE-TUR
\]

Considerando o govenrno de forma explícita vamos ter a Rena Líquida do Governo (RLG) composta pela soma dos impostos diretos e indiretos e das outras receitas correntes do governo, deduzidas as transferências. Então podemos reescrever:

\[
RDB=RPD+RLG
\]

Onde RPD é a Renda Privada Disponível. Considerando a naureza da demanda e também a alocação da renda gerada, retomemos por pouco tempo a ideia de uma economia fechada e sem governo. Considerando tal economia, teríamos então por um lado, o produto (Y) é igual à demanda por consumo (C) e à demanda por investimento (I), enquanto, por outro, a renda (Y) é igual a consumo (C) mais poupança (S). Formalmente temos que $Y=C+I$ e $Y=C+S$, então $I=S$.

Introduzindo o governo, temos uma demanda adicional através dos gastos do governo $G$, então $Y=C+I+G$, e também uma renda adicional atraves da Renda Líquida do Governo RLG, então $Y=C+S+RLG$ ou 

\[
I=S+RLG-G=S+S_{g}=SD
\]

Onde $S_{g}=RLG-G$ é a poupança do governo, e agora $S$ é explicitamente a popupança privada. Também denotamos $SD=S+S_{g}$ a poupança doméstica, isto é, a poupança bruta da economia. Isso posto, temos por fim de reintroduzir o setor externo na identidade. 

\[
I=S+S_{g}+S_{e}
\]
Que é definida como $s_{e}=\left(M-X\right)+RLEE-TUR$. Podemos demonstrar que podemos obter isso setemos $Y=C+I+G+\left(X-M\right)=PIB$ ao olharmos a natureza da demanda, e $Y=C+S+RLG+RLEE-TUR=RDB+RLEE-TUR=PIB$ se olharmos a alocação da renda\footnote{Note que $RDB=C+S+RLG$, isto é, a Renda Disponível Bruta é a soma entre o consumo e poupança privados, e a renda líquida do governo. E como tínhamos ainda que $RDB=RPD+RLG$ então $RPD=C+S$, ou seja, a renda privada disponível é a soma do consumo e poupança.}. Então:

\begin{align*}
I+G+\left(X-M\right) & =S+RLG+RLEE-TUR\\
I & =S+\left[RLG-G\right]+\left[\left(M-X\right)+RLEE-TUR\right]\\
 & =S+S_{g}+S_{e}
\end{align*}

Lembrando que todas essas discussões usa o SNA 93, que não é o mais atual, porém, eu acredito que seria interessante do ponto de vista didático para entender algumas relações matemáticas, e também porque pode ser útil para futuras modelagens.

\subsubsection{Problemas de mensuração}

As dificuldades envolvidas na mensuração dos agregados são de três tipos: 
\begin{enumerate}
\item Dificuldades técnicas: as principais dificuldades técnicas deste tipo são
\begin{enumerate}
\item a existência de variação nos preços em momentos distintos do tempo e, com isso a necessidade de se distinguir entre contabilidade real e contabilidade nominal; 
\item a necessidade de se estabelecer comparações entre países em que os agregados são originalmente estimados em moedas distintas, o que gerou o conceito de paridade de poder de compra (PPP, na sigla em inglês purchase power parity); 
\end{enumerate}
\item Dificuldades operacionais: estão relacionadas à existência de atividades econômicas que não são efetuadas por meio de empresas formalmente constituídas e/ou por empresas formalmente constituídas, mas que desempenham parte de suas atividades de modo informal.
\item Dificuldades conceituais: relacionam-se à existência de atividades não monetizadas e aos problemas com o meio ambiente.
\end{enumerate}
Quando se compara um agregado, por exemplo o PIB, em diferentes momentos do tempo, para verificar sua evolução, é preciso se certificar de que a comparação é válida. Do contrário, poderemos ser induzidos a achar que, por exemplo, houve crescimento de uma dada economia em um determinado período, quando, na realidade, pode ter havido apenas elevação dos preços, não pum crescimento real. Para que a comparação seja válida, é preciso comparar os agregados mensurando-os com preços do mesmo momento do tempo. Sem tomar esse cuidado, estaremos observando apenas as variações nominais. A técnica que nos permite comparar o valor das variáveis econômicas a preços de um mesmo momento chama-se deflacionamento e seus instrumentos são os índices de preço\footnote{Um índice de preços é uma medida estatística que sintetiza a variação dos preços de um conjunto de bens e serviços ao longo do tempo em relação a um período de referência. Sua principal função é mensurar a inflação ou deflação e permitir a comparação de valores monetários observados em momentos distintos por meio do deflacionamento.}.

Estes índices de preço mensuram a variação dos preços em um determinado período de tempo e aparecem sob a forma de números-índice\footnote{Número-índice é uma medida estatística que expressa a variação relativa de uma variável em relação a um período-base, geralmente fixado em 100, permitindo comparar sua evolução ao longo do tempo.}. Os índices simples buscam medir a evolução de apenas uma série homogênea de dados, por exemploo se olhamos a produção de trigo a partir dos anos 2000, podemos definir a produção como 100 neste ano, e reescrever a produção nos próximos anos como reescalado pelo valor do ano 2000, por exemplo se em 2001 produção dobrou, então a produção agora é 200. 

Já os índices compostos são utilizados quando se torna necessário trabalhar com um conjunto de séries de natureza distinta, por exemplo, preços e quantidades. 

Os mais conhecidos dentre os índices de preço são os de Laspeyres\footnote{Mantém fixa as quantidades do período-base: Quanto custaria hoje a cesta que era consumida no ano-base?}, Paasche\footnote{Mantém fixa as quantidades do período atual: Quanto custaria a cesta atual se fosse comprada aos preços antigos?} e Fisher\footnote{Combina os dois, uma média geométrica, isto é $\sqrt{AB}$}.

Há também uma forma indireta de se mensurar a evolução dos preços: o deflator implícito do PIB. Trata-se de uma forma indireta porque o índice é construído não a partir do monitoramento dos preços, mas a partir da comparação entre os valores a preços correntes estimados para o produto agregado e os índices de crescimento real apurados a cada período.

A contabilidade nacional distingue entre juros nominais e juros reais porque a inflação incide diretamente sobre o valor dos ativos financeiros de valor nominal constante. Os ajustes contábeis derivados da existência de inflação em um determinado ano incidem apenas sobre a distribuição da renda entre os diferentes agentes e não sobre o montante dos agregados.

12. Em um contexto inflacionário, a renda nominal de um agente qualquer tenderá a ser maior do que a real se ele for liquidamente um credor, ao passo que a situação deverá ser inversa se ele for liquidamente um devedor.

13. A existência hoje de intensos fluxos de capital entre as diferentes economias e o fato de que a maior parte dos países adota um regime de câmbio em que a taxa de câmbio é determinada pelo mercado (regime de câmbio flutuante) torna essa variável muito mais um reflexo da relação entre oferta e demanda por divisas no contexto do mercado de capitais, do que uma expressão do efetivo poder de compra das diferentes moedas aí envolvidas.

14. Além disso, nem todos os bens e serviços produzidos são tradables (negociados internacionalmente) e há diferenças na política tarifária e de subsídios, e nos custos de transporte entre as diferentes economias. Todos esses fatores fazem com que a taxa de câmbio não seja um conversor eficiente para se comparar o produto de diferentes países A taxa de câmbio PPP (purchase power parity) procura resolver esse problema.

15. O conhecimento sobre a forma como a economia está se comportando é fundamental para a tomada de decisões sobre produção, investimento e consumo, tanto no âmbito privado quanto público.

17. A incorporação às estimativas do produto agregado do valor produzido pela chamada economia informal enfrenta dificuldades de natureza operacional, uma vez que é bastante difícil identificar e localizar as atividades que a constituem. 

8. A economia subterrânea, desenvolvida por empresas formais ou informais, na qual as atividades não são declaradas com o deliberado intuito de escapar de tributos, leis trabalhistas e outras normas legais, é parte da economia informal. Além dessa parte, estimada hoje no Brasil pelo Índice da Economia Subterrânea, a economia informal envolve a atividade produtiva das famílias exercida sem a constituição formal de empresas e as atividades ilegais. Parte da atividade produtiva informal das famílias é capturada pelas CEI Institucionais.

19. A contabilidade nacional procura estimar o valor monetário das atividades não monetizadas, imputando-lhes os valores que elas supostamente teriam se tivessem passado pelo mercado. A decisão sobre quais dessas atividades devem integrar o SCN é quase sempre convencional, uma vez que conceitualmente não há consenso sobre o tema.

20. A expansão acelerada e sem controle da industrialização e das formas urbanas de vida tem provocado a degradação do meio ambiente e tem esbarrado nos limites impostos pelo estoque finito de recursos naturais do planeta.

21. As pressões sobre o meio ambiente decorrentes da produção e do consumo constituem externalidades negativas, ou seja, custos não valorados pelo mercado e que deveriam de alguma maneira afetar os valores apresentados pelas contas nacionais.

22. Apesar das dificuldades teóricas e técnicas envolvidas nessas estimativas, organismos como a ONU e a OCDE vêm investindo na criação de IDS (Indicadores de Desenvolvimento Sustentável), os quais auxiliariam na tarefa de proceder essa mensuração. No Brasil o IBGE também calcula os IDS. Em sua última versão, de 2004, são 54 indicadores, organizados em quatro dimensões, tal como recomendado pela ONU (social, ambiental, econômica e institucional).

\bibliographystyle{plain}
\bibliography{tese,Artigo1/references,Artigo2/referencias,Artigo3/referencias}

\end{document}
